\section{Trace, endpoint jet, and restart}
\label{sec:downstream}

The datum-generated Theorem~\ref{thm:producer} supplies
Proposition~\ref{prop:projective} and \eqref{eq:pi0} on the
entire alleged finite maximal interval \(I_*\).

\subsection{The pressure-complete equation in every spatial row}

Let \(R_j=\partial_j(-\Delta)^{-1/2}\) be the Riesz transforms and let
\(\Leray\) be the Leray projection. With the decaying normalization,
\begin{equation}
 p=R_iR_j(u_i u_j),
 \qquad
 (u\cdot\nabla)u+\nabla p
 =\Leray\nabla\cdot(u\otimes u).
 \label{eq:pressure-projection}
\end{equation}
The first identity follows by taking the divergence of \eqref{eq:NS}; the
second is the same equation written on the divergence-free subspace. Pressure
has not been discarded: it is the simultaneous whole-space response in
\eqref{eq:pressure-projection}.

\begin{lemma}[The complete first spatial row reaches every endpoint topology]
\label{lem:S0-endpoint}
If \(u\in\Szero(I_*)\), then, for every integer \(m\geq0\),
\[
 u_t\in L^1(I_*;H^m),
 \qquad
 u\in W^{1,1}(I_*;H^m)\hookrightarrow C([0,T_*];H^m).
\]
\end{lemma}

\begin{proof}
Because \(H^{m+2}(\R^3)\) is an algebra and the Riesz transforms and
\(\Leray\) are bounded Fourier multipliers of order zero,
\begin{equation}
 \norm{\Leray\nabla\cdot(u\otimes u)}_{H^m}
 \leq C_m\norm{u}_{H^{m+2}}^2.
 \label{eq:nonlinear-Hm}
\end{equation}
Consequently,
\begin{align}
 \norm{\Leray\nabla\cdot(u\otimes u)}_{L^1(I_*;H^m)}
 &\leq C_m\norm{u}_{L^2(I_*;H^{m+2})}^2<\infty,
 \label{eq:nonlinear-L1}\\
 \norm{\Delta u}_{L^1(I_*;H^m)}
 &\leq T_*^{1/2}\norm{u}_{L^2(I_*;H^{m+2})}<\infty.
 \label{eq:viscous-L1}
\end{align}
Equation~\eqref{eq:NS}, equivalently
\[
 u_t=\nu\Delta u-\Leray\nabla\cdot(u\otimes u),
\]
then gives \(u_t\in L^1(I_*;H^m)\). Since
\(u\in L^2(I_*;H^m)\subset L^1(I_*;H^m)\), the Bochner fundamental theorem
of calculus gives a unique representative in \(C([0,T_*];H^m)\).
\end{proof}

The full mixed intersection also contains the sharper adjacent coordinates
that explain the trace mechanism directly.

\begin{lemma}[Adjacent Hilbert-scale trace]
\label{lem:LM}
Let \(s\in\R\). If
\[
 v\in L^2(I_*;H^{s+1}),
 \qquad v_t\in L^2(I_*;H^{s-1}),
\]
then \(v\) has a unique representative in \(C([0,T_*];H^s)\). For
\(0\leq a\leq b\leq T_*\), this representative satisfies
\begin{equation}
 \left|\norm{v(b)}_{H^s}^2-\norm{v(a)}_{H^s}^2\right|
 \leq
 2\norm{v_t}_{L^2(a,b;H^{s-1})}
  \norm{v}_{L^2(a,b;H^{s+1})}.
 \label{eq:LM-estimate}
\end{equation}
\end{lemma}

\begin{proof}
The Hilbert scale
\[
 H^{s+1}\hookrightarrow H^s\hookrightarrow H^{s-1}
\]
is a Gelfand triple with pivot \(H^s\); this is the standard
Lions--Magenes trace setting \cite{LM}. Time mollification gives, first for
smooth functions and then by density,
\[
 \norm{v(b)}_{H^s}^2-\norm{v(a)}_{H^s}^2
 =2\int_a^b
 \pair{v_t(t)}{v(t)}_{H^{s-1},H^{s+1};H^s}\,dt.
\]
The shifted dual pairing obeys
\[
 \left|\pair{f}{g}_{H^{s-1},H^{s+1};H^s}\right|
 \leq\norm{f}_{H^{s-1}}\norm{g}_{H^{s+1}},
\]
so Cauchy--Schwarz proves \eqref{eq:LM-estimate}. The same approximation
argument gives weak \(H^s\)-continuity; weak continuity together with the
continuity of the squared norm supplied by the identity gives strong
\(H^s\)-continuity. Uniqueness follows because two continuous representatives
agreeing almost everywhere agree everywhere.
\end{proof}

Applying Lemma~\ref{lem:LM} to every temporal row gives, for all
\(N,m\in\N_0\),
\begin{equation}
 \partial_t^N u\in C([0,T_*];H^m).
 \label{eq:all-traces}
\end{equation}
Indeed, the two hypotheses are precisely adjacent coordinates of
\(\cI[u_0]\). This is the Lions--Magenes trace mechanism behind the endpoint,
whereas Lemma~\ref{lem:S0-endpoint} shows that the projected row
\(\Szero(I_*)\) already reaches the same endpoint through the equation.

\subsection{The adjacent critical block}

Put \(\Lambda=(-\Delta)^{1/2}\) and
\begin{equation}
 E_s(t):=\frac12\norm{\Lambda^s u(t)}_2^2,
 \qquad H(t):=E_{1/2}(t).
 \label{eq:energies}
\end{equation}

\begin{proposition}[Critical trace and current]
\label{prop:critical}
The adjacent coordinates in \(\cI[u_0]\) imply
\begin{align}
 u\in L^2(I_*;H^{5/2}),\quad
 u_t\in L^2(I_*;H^{1/2})
 &\Longrightarrow
 E_{3/2}\in L^\infty(I_*)\cap L^1(I_*),
 \label{eq:critical-first}\\
 u,u_t\in L^2(I_*;H^1)
 &\Longrightarrow H\in W^{1,1}(I_*).
 \label{eq:critical-second}
\end{align}
Moreover the pressure-explicit critical current
\begin{equation}
 \mathcal P_{1/2}
 :=-\pair{\Lambda^{1/2}u}
 {\Lambda^{1/2}\bigl((u\cdot\nabla)u+\nabla p\bigr)}
 \label{eq:critical-current}
\end{equation}
satisfies
\begin{equation}
 \mathcal P_{1/2}=H'+2\nu E_{3/2}\in L^1(I_*).
 \label{eq:critical-balance}
\end{equation}
\end{proposition}

\begin{proof}
Lemma~\ref{lem:LM} with \(s=3/2\) gives
\(u\in C([0,T_*];H^{3/2})\). More precisely, the same
time-mollification argument gives, for \(0\leq r\leq t\leq T_*\),
\begin{equation}
 E_{3/2}(t)-E_{3/2}(r)
 =\int_r^t\operatorname{Re}
 \bigl(\Lambda^{1/2}u_t(s),\Lambda^{5/2}u(s)\bigr)_2\,ds.
 \label{eq:critical-trace-identity}
\end{equation}
Cauchy--Schwarz therefore gives \(E_{3/2}\in W^{1,1}(I_*)\), and in
particular \(E_{3/2}\in L^\infty(I_*)\). Also
\[
 \int_0^{T_*}E_{3/2}(t)\,dt
 \leq\frac12\norm{u}_{L^2(I_*;H^{5/2})}^2<\infty,
\]
which proves \eqref{eq:critical-first}.

The second pair means \(u\in H^1(I_*;H^1)\). The Hilbert-space chain rule
therefore yields
\begin{equation}
 H'(t)=\pair{\Lambda^{1/2}u(t)}{\Lambda^{1/2}u_t(t)}
 \quad\text{for a.e. }t,
 \label{eq:H-prime}
\end{equation}
and
\[
 \norm{H'}_{L^1(I_*)}
 \leq\norm{u}_{L^2(I_*;H^1)}
       \norm{u_t}_{L^2(I_*;H^1)}<\infty.
\]
Thus \(H\in W^{1,1}(I_*)\).

On every compact subinterval of \(I_*\), pair \eqref{eq:NS} with
\(\Lambda u\). The viscous term is
\(-\nu\norm{\Lambda^{3/2}u}_2^2=-2\nu E_{3/2}\), so
\[
 H'=-2\nu E_{3/2}+\mathcal P_{1/2}.
\]
The pressure in \eqref{eq:critical-current} is the normalized solution of the
Poisson equation in \eqref{eq:pressure-projection}; equivalently, its global
pairing vanishes because \(\nabla\cdot\Lambda u=0\). The pressure-explicit
form is retained because it is the form differentiated in the jet below.
Since \(H'\) and \(E_{3/2}\) are both in \(L^1(I_*)\), the classical current
extends to the whole terminal interval as the \(L^1\) function in
\eqref{eq:critical-balance}.
\end{proof}

No additional critical-current estimate is required after the compatible
intersection is present: Proposition~\ref{prop:critical} is a trace consequence
of two adjacent rectangles.

\subsection{The unique projective endpoint}

For each \(m\), let \(u_*^{(m)}\in H^m\) denote the endpoint value furnished
by Lemma~\ref{lem:S0-endpoint}, equivalently by Lemma~\ref{lem:LM}. If
\(k>m\), then the \(H^k\)-continuous and \(H^m\)-continuous representatives
are the same function almost everywhere on \(I_*\); after embedding
\(H^k\hookrightarrow H^m\), continuity makes them identical on
\([0,T_*]\). In particular their endpoint values agree in \(H^m\).
Therefore the family \(\{u_*^{(m)}\}_m\) is compatible and defines one element
\begin{equation}
 u_*:=u(T_*)\in\bigcap_{m=0}^{\infty}H^m(\R^3)=\Hs(\R^3).
 \label{eq:endpoint}
\end{equation}
No bound uniform in \(m\) is asserted or needed. Since divergence is
continuous \(H^m\to H^{m-1}\), \(u_*\) is divergence-free; since the
\(H^0\)-trace is strong, it also has finite kinetic energy.

The same argument applied to \(\partial_t^N u\) in \eqref{eq:all-traces}
gives compatible endpoint traces
\begin{equation}
 V_N^*:=\lim_{t\uparrow T_*}\partial_t^Nu(t)\in\Hs,
 \qquad N\in\N_0.
 \label{eq:velocity-jet-traces}
\end{equation}

\subsection{The pressure-explicit endpoint jet}

For a multi-index \(\alpha\), set
\[
 J_{N,\alpha}:=\partial_t^N\partial_x^\alpha u,
 \qquad
 \Pi_{N,\alpha}:=\partial_t^N\partial_x^\alpha p.
\]
Direct differentiation of \eqref{eq:NS} gives, on \(I_*\),
\begin{equation}
\begin{aligned}
 J_{N+1,\alpha}
 &+
 \sum_{a=0}^{N}\sum_{\beta\leq\alpha}
 \binom Na\binom\alpha\beta
 (J_{a,\beta}\cdot\nabla)
 J_{N-a,\alpha-\beta}
 +\nabla\Pi_{N,\alpha}
 \\
 &=\nu\Delta J_{N,\alpha},
 \qquad \nabla\cdot J_{N,\alpha}=0,
\end{aligned}
\label{eq:full-binomial}
\end{equation}
while incompressibility fixes every pressure row through
\begin{equation}
 -\Delta\Pi_{N,\alpha}
 =\partial_i\partial_j\partial_t^N\partial_x^\alpha(u_i u_j).
 \label{eq:pressure-poisson-jet}
\end{equation}

For the pure-time jet write \(V_N=\partial_t^Nu\) and
\(Q_N=\partial_t^Np\). Equations
\eqref{eq:full-binomial}--\eqref{eq:pressure-poisson-jet} become
\begin{align}
 Q_N
 &=R_iR_j\sum_{a=0}^{N}\binom Na V_{a,i}V_{N-a,j},
 \label{eq:pressure-recurrence}\\
 V_{N+1}
 &=\nu\Delta V_N
 -\sum_{a=0}^{N}\binom Na(V_a\cdot\nabla)V_{N-a}
 -\nabla Q_N.
 \label{eq:velocity-recurrence}
\end{align}

\begin{proposition}[Endpoint jet reconstruction]
\label{prop:endpoint-jet}
The endpoint velocity \(u_*\) uniquely determines a sequence
\((V_N^*,Q_N^*)_{N\geq0}\subset\Hs\times\Hs\) by
\begin{align}
 V_0^*&=u_*,
 \label{eq:jet-start}\\
 Q_N^*&=R_iR_j\sum_{a=0}^{N}\binom Na V_{a,i}^*V_{N-a,j}^*,
 \label{eq:jet-pressure}\\
 V_{N+1}^*&=\nu\Delta V_N^*
 -\sum_{a=0}^{N}\binom Na(V_a^*\cdot\nabla)V_{N-a}^*
 -\nabla Q_N^*.
 \label{eq:jet-velocity}
\end{align}
This recursively constructed sequence equals the left endpoint traces of the
actual velocity-pressure jet.
\end{proposition}

\begin{proof}
The Fr\'echet space \(\Hs\) is closed under spatial differentiation and
pointwise multiplication. Riesz transforms preserve each \(H^m\), hence
preserve \(\Hs\). Thus \eqref{eq:jet-pressure}--\eqref{eq:jet-velocity}
construct one and only one \(\Hs\) pair at every order.

It remains to identify this formal recursion with the history. Fix \(m\).
The traces \eqref{eq:velocity-jet-traces} converge in arbitrarily high
Sobolev spaces. Sobolev multiplication, spatial differentiation, and the
Riesz transforms are continuous between the corresponding spaces. We may
therefore pass \(t\uparrow T_*\) in
\eqref{eq:pressure-recurrence}--\eqref{eq:velocity-recurrence}. Beginning
with \(V_0^*=u_*\), induction in \(N\) gives exactly
\eqref{eq:jet-pressure}--\eqref{eq:jet-velocity}. Since \(m\) was arbitrary,
the identities hold in \(\Hs\). Spatial differentiation then identifies the
full multi-index jet in \eqref{eq:full-binomial}.
\end{proof}

\subsection{Classical restart and contradiction of finite maximality}

We record the standard local theorem in precisely the form used here
\cite{Temam}.

\begin{lemma}[Smooth local restart and uniqueness]
\label{lem:local}
Let \(t_0\in\R\) and let \(a\in\Hs_\sigma(\R^3)\).  There is a
\(\delta>0\) and a unique smooth Navier--Stokes solution \((w,q)\) on
\([t_0,t_0+\delta]\times\R^3\) with \(w(t_0)=a\).  Every spatial Sobolev
order persists on this common interval, and the temporal jet at \(t_0\) is
the pressure-normalized recursion
\eqref{eq:jet-pressure}--\eqref{eq:jet-velocity} with \(V_0^*=a\).
\end{lemma}

\begin{proof}
This is the classical smooth-data local theorem for three-dimensional
Navier--Stokes \cite{Temam}.  Applied to \(a\in H^\infty\), it gives one
positive lifespan on which every higher Sobolev norm propagates along the same
solution.  Smoothness up to the initial face follows from the equation.
Differentiating there and using the normalized pressure Poisson equation gives
the stated jet recursion.  Uniqueness is the standard strong-solution
uniqueness on the same interval.
\end{proof}

Apply Lemma~\ref{lem:local} at time \(T_*\) with datum \(a=u_*\). It gives a
smooth solution \((w,q)\) on \([T_*,T_*+\delta]\). Define
\[
 (\widetilde u,\widetilde p)(t)
 :=
 \begin{cases}
 (u,p)(t),&0\leq t<T_*,\\
 (w,q)(t),&T_*\leq t\leq T_*+\delta.
 \end{cases}
\]
The velocity traces agree at \(T_*\) by construction. Therefore the
piecewise function has no distributional time-boundary defect and solves
\eqref{eq:NS} across \(T_*\). Proposition~\ref{prop:endpoint-jet} and the
last assertion of Lemma~\ref{lem:local} show that every left and right
velocity-pressure derivative agrees at \(T_*\). Hence
\((\widetilde u,\widetilde p)\) is classical, indeed smooth, on an interval
strictly larger than \([0,T_*)\).

Equivalently, one may start the local solution at a time
\(t_0<T_*\) sufficiently close to \(T_*\).  Continuity of the endpoint in
the fixed smooth-data local topology gives a common positive local lifespan,
and uniqueness identifies that solution with \(u\) on the overlap.  Both
descriptions restart the same datum-generated history.

This contradicts the definition of \(T_*\) as the maximal classical
lifespan. Thus the assumption \(T_*<\infty\) is impossible, and
\[
 \boxed{T_*=\infty.}
\]
Because the argument applies at every finite time and all velocity and
pressure rows lie in \(\Hs\), Sobolev embedding yields
\[
 \boxed{u,p\in C^\infty([0,\infty)\times\R^3).}
\]
This completes every implication in \eqref{eq:chain} and proves
Theorem~\ref{thm:main}.
